% Gerado por scripts/migrate_markdown_to_latex.py.
% Fonte migrada: fases/01_validacao_conceitual/docs/confronto_resultados_literatura.md


\chapter{Confronto entre resultados observados e literatura}





Este documento responde uma pergunta central da validacao: os numeros obtidos nos testes batem com o comportamento esperado pelas definicoes aceitas na literatura?





A resposta curta e: sim para os testes matematicos e de aderencia conceitual.
A bateria fisica posterior confirmou representacoes e kernels INT8/2:4, mas
contradisse ganho de latencia nas operacoes isoladas. A comparacao de desempenho
com modelos completos e setups equivalentes aos artigos continua em andamento.



\section{Criterio usado}





Os artigos e documentacoes primarias nem sempre fornecem os mesmos tensores pequenos que usamos nos testes. Por isso, a comparacao inicial nao e 'mesmo dataset, mesmo resultado final do artigo'. A comparacao correta nesta fase e:


\begin{enumerate}
\item pegar a definicao publicada;
\item derivar um caso pequeno e deterministico;
\item calcular o valor esperado pela formula;
\item comparar a implementacao local com esse valor;
\item limitar a conclusao ao nivel que o teste realmente sustenta.
\end{enumerate}


\section{Resultados numericos verificados}



\begin{landscape}
\scriptsize
\begin{longtable}{@{}>{\RaggedRight\arraybackslash}p{0.137\linewidth}>{\RaggedRight\arraybackslash}p{0.137\linewidth}>{\RaggedRight\arraybackslash}p{0.137\linewidth}>{\RaggedRight\arraybackslash}p{0.137\linewidth}>{\RaggedRight\arraybackslash}p{0.137\linewidth}>{\RaggedRight\arraybackslash}p{0.137\linewidth}@{}}
\toprule
Conceito & Fonte usada & Valor esperado pela definicao & Resultado observado & Bate? & Limite da conclusao \\
\midrule
\endfirsthead
\toprule
Conceito & Fonte usada & Valor esperado pela definicao & Resultado observado & Bate? & Limite da conclusao \\
\midrule
\endhead
Projecao densa & Goodfellow et al.; Vaswani et al. (2017) para projecoes lineares em Q/K/V & Para \texttt{X=[[1,2],[3,4]]}, \texttt{W=[[1,0],[0,2]]}, \texttt{b=[0.\allowbreak{}5,-0.\allowbreak{}5]}, a saida esperada e \texttt{[[1.\allowbreak{}5,3.\allowbreak{}5],[3.\allowbreak{}5,7.\allowbreak{}5]]}. & O teste retornou exatamente \texttt{[[1.\allowbreak{}5, 3.\allowbreak{}5], [3.\allowbreak{}5, 7.\allowbreak{}5]]}. & Sim & Valida a operacao linear \texttt{Y=XW\textasciicircum{}T+b}, nao uma rede completa. \\
Scaled dot-product attention & Vaswani et al. (2017), formula \texttt{softmax(QK\textasciicircum{}T/\allowbreak{}sqrt(d\_\allowbreak{}k))V} & Para \texttt{Q=K=[[1,0],[0,1]]}, \texttt{V=[[1,2],[3,4]]}, a saida esperada e \texttt{[[1.\allowbreak{}6604769, 2.\allowbreak{}6604769], [2.\allowbreak{}3395231, 3.\allowbreak{}3395231]]}. & O teste retornou \texttt{[[1.\allowbreak{}6604769, 2.\allowbreak{}6604769], [2.\allowbreak{}3395228, 3.\allowbreak{}3395233]]}. & Sim & Valida a operacao matematica de atencao escalonada, nao uma arquitetura Transformer completa. \\
Attention PyTorch & Documentacao PyTorch + Vaswani et al. (2017) & A API deve produzir a mesma operacao de scaled dot-product attention, com tolerancia numerica. & A saida de \texttt{torch.\allowbreak{}nn.\allowbreak{}functional.\allowbreak{}scaled\_\allowbreak{}dot\_\allowbreak{}product\_\allowbreak{}attention} bateu com a implementacao manual em \texttt{1e-6}. & Sim & Valida a API como referencia algoritmica para attention pequena e deterministica. \\
Attention Keras/TensorFlow & Documentacao Keras Ops + Vaswani et al. (2017) & A API publica deve produzir a mesma scaled dot-product attention para os mesmos \texttt{Q}, \texttt{K}, \texttt{V} e mascara. & Os tres casos contra NumPy foram aprovados; o maior erro absoluto Keras/TensorFlow foi \texttt{2.\allowbreak{}38418579e-07}, abaixo de \texttt{atol=1e-6}. & Sim & Valida equivalencia numerica entre frameworks em CPU/float32; nao compara desempenho. \\
Mascara em attention & Documentacao PyTorch SDPA & Mascara aditiva deve alterar os scores antes do \texttt{softmax} e produzir a mesma saida da API de referencia. & A implementacao manual com mascara bateu com PyTorch em \texttt{1e-6}. & Sim & Valida mascara aditiva no caso pequeno; outros tipos de mascara devem ser testados se usados. \\
Multi-head attention & Vaswani et al. (2017), divisao da representacao em multiplas heads & A saida apos split, attention por head e combine deve preservar shape \texttt{[batch, seq, d\_\allowbreak{}model]}. & Para entrada \texttt{[2,3,4]} e \texttt{2} heads, a saida preservou \texttt{[2,3,4]} e bateu com PyTorch em \texttt{1e-6}. & Sim & Valida a mecanica multi-head, ainda sem treino ou avaliacao de tarefa. \\
Dependencia entre posicoes & Vaswani et al. (2017), ideia de self-attention sobre sequencias & Alterar o valor de uma posicao pode alterar a saida calculada em outra posicao. & O teste alterou um token e a saida de outra posicao mudou. & Sim & Valida dependencia sequencial minima; ainda nao prova qualidade de modelagem. \\
Bloco Transformer simplificado & Vaswani et al. (2017) + criterios do documento teorico & Deve conter entrada sequencial, Q/K/V, attention escalonada, posicao, residual, norm e FFN. & O bloco preservou shape \texttt{[2,3,4]} e foi classificado como \texttt{bloco\_\allowbreak{}transformer\_\allowbreak{}simplificado}. & Sim & Valida um bloco controlado; nao deve ser chamado de Transformer completo. \\
Auditoria Transformer da NN atual & Vaswani et al. (2017) + protocolo do projeto & Um bloco Transformer deve passar em Q/K/V, attention escalonada, dependencia posicional, informacao posicional, residual, norm e FFN; uma NN comum nao deve passar. & A NN atual passou como \texttt{bloco\_\allowbreak{}transformer\_\allowbreak{}simplificado}; uma NN linear comum foi rejeitada como \texttt{nao\_\allowbreak{}aderente}. & Sim & Valida que o nucleo pertence a familia Transformer, mas nao autoriza chama-lo de Transformer completo. \\
KV cache & Literatura de inferencia autoregressiva e compressao de KV cache & Saida com cache deve bater com saida causal completa e reduzir recomputacao de K/V. & Para \texttt{4} tokens, cache projetou \texttt{4} K/V contra \texttt{10} no caminho ingenuo; saidas bateram em \texttt{1e-6}. & Sim & Valida equivalencia e reaproveitamento algoritmico; ganho de hardware ainda exige medicao. \\
FLOPs teoricos & Hennessy e Patterson; analise de complexidade de atencao em Vaswani et al. (2017) & Para o caso pequeno definido, projecao densa \texttt{270}, attention \texttt{288}, bloco \texttt{1992}. & Os testes retornaram \texttt{270}, \texttt{288} e \texttt{1992}. & Sim & Valida contagem teorica do recorte; FLOPs nao equivalem automaticamente a latencia. \\
MSE & Definicao estatistica e \texttt{sklearn.\allowbreak{}metrics.\allowbreak{}mean\_\allowbreak{}squared\_\allowbreak{}error} & Para \texttt{y=[1,2,3]}, \texttt{y\_\allowbreak{}hat=[1,2,4]}, \texttt{MSE=(0\textasciicircum{}2+0\textasciicircum{}2+1\textasciicircum{}2)/\allowbreak{}3=0.\allowbreak{}3333333}. & \texttt{0.\allowbreak{}3333333333333333}. & Sim & Valida a metrica numerica. \\
MAE & Definicao estatistica e \texttt{sklearn.\allowbreak{}metrics.\allowbreak{}mean\_\allowbreak{}absolute\_\allowbreak{}error} & Para o mesmo exemplo, \texttt{MAE=(0+0+1)/\allowbreak{}3=0.\allowbreak{}3333333}. & \texttt{0.\allowbreak{}3333333333333333}. & Sim & Valida a metrica numerica. \\
R2 & Definicao de coeficiente de determinacao e \texttt{sklearn.\allowbreak{}metrics.\allowbreak{}r2\_\allowbreak{}score} & Para o mesmo exemplo, \texttt{R2=1-SS\_\allowbreak{}res/\allowbreak{}SS\_\allowbreak{}tot=1-1/\allowbreak{}2=0.\allowbreak{}5}. & \texttt{0.\allowbreak{}5}. & Sim & Valida a metrica numerica. \\
Similaridade de cosseno & Manning, Raghavan e Schutze; scikit-learn pairwise metrics & Para \texttt{[1,1,0]} e \texttt{[1,0,0]}, \texttt{cos=1/\allowbreak{}sqrt(2)=0.\allowbreak{}70710678}. & \texttt{0.\allowbreak{}7071067811865475}. & Sim & Valida a metrica numerica. \\
Pruning L1 PyTorch & Tutorial/API PyTorch pruning & Em matriz \texttt{2x4}, \texttt{amount=0.\allowbreak{}5} deve zerar 4 de 8 pesos e criar mascara. & \texttt{4/\allowbreak{}8} zeros, esparsidade \texttt{0.\allowbreak{}5}, mascara \texttt{weight\_\allowbreak{}mask} presente. & Sim, parcialmente & Valida pruning conceitual/algoritmico; nao valida ganho real de hardware porque o caminho continua denso. \\
Quantizacao simetrica int8 manual & Literatura de compressao/quantizacao e definicao operacional & Para tensor de 5 valores \texttt{float32}, esperado reduzir de \texttt{20} bytes para \texttt{5} bytes em armazenamento int8; escala \texttt{2/\allowbreak{}127=0.\allowbreak{}015748}. & \texttt{dtype=torch.\allowbreak{}int8}, \texttt{5} bytes contra \texttt{20} bytes; erro maximo \texttt{0.\allowbreak{}007874}, dentro do limite de arredondamento da escala. & Sim & Valida quantizacao numerica; nao valida kernel de inferencia. \\
torchao int8 weight-only v2 & Documentacao torchao \texttt{Int8WeightOnlyConfig} & Configuracao deve aplicar quantizacao simetrica int8 weight-only em camadas lineares. & Peso \texttt{Int8Tensor} possui \texttt{qdata} em \texttt{int8}; o profiler fisico confirmou \texttt{weight\_int8pack\_mm}. & Sim & Valida representacao compacta e caminho weight-only; nao implica ativacoes inteiras. \\
\bottomrule
\end{longtable}
\end{landscape}



\section{O que ainda nao esta validado}






Ja podemos relatar compressao, memoria e latencia das operacoes isoladas na RTX
4070 Ti SUPER. Ainda nao podemos afirmar reproducao de artigos de desempenho:
os setups, modelos e kernels publicados nao sao necessariamente equivalentes, e
a avaliacao final dos tres modelos OPT permanece pendente.



Para confrontar desempenho diretamente com a literatura, o protocolo precisa fixar:


\begin{enumerate}
\item modelo ou bloco usado;
\item tamanho de sequencia;
\item batch size;
\item dtype;
\item hardware;
\item kernel usado;
\item metrica medida;
\item baseline;
\item tolerancia de erro numerico;
\item repeticoes e intervalo de confianca.
\end{enumerate}


\section{Conclusao metodologica}






As ferramentas atuais sao adequadas para iniciar o trabalho porque passam na validacao de definicao: projecao densa, attention, bloco simplificado, KV cache, FLOPs, metricas, pruning conceitual e quantizacao numerica se comportam como esperado pelas formulas e documentacoes primarias.






Os caminhos fisicos agora sao suficientes para afirmar reducao de armazenamento
e memoria no recorte medido, pois formato e kernels foram confirmados. Eles nao
sustentam ganho de latencia: os seis grupos ficaram abaixo do baseline. Uma
generalizacao para Transformers pre-treinados depende da coleta OPT e uma
comparacao direta com artigos exige equivalencia de setup.



\section{Referencias}


\begin{itemize}
\item Vaswani et al. (2017). \emph{Attention Is All You Need}:
.
\item PyTorch. \texttt{scaled\_\allowbreak{}dot\_\allowbreak{}product\_\allowbreak{}attention}:
.
\item Keras. \texttt{keras.\allowbreak{}ops.\allowbreak{}nn.\allowbreak{}dot\_\allowbreak{}product\_\allowbreak{}attention}:
.
\item PyTorch. \texttt{torch.\allowbreak{}nn.\allowbreak{}utils.\allowbreak{}prune.\allowbreak{}l1\_\allowbreak{}unstructured}:
.
\item PyTorch Tutorials. \emph{Pruning Tutorial}:
.
\item PyTorch. \texttt{torchao.\allowbreak{}quantization.\allowbreak{}Int8WeightOnlyConfig}:
.
\item scikit-learn. \texttt{mean\_\allowbreak{}squared\_\allowbreak{}error}:
.
\item scikit-learn. \texttt{r2\_\allowbreak{}score}:
.
\item scikit-learn. \emph{Cosine similarity}:
.
\end{itemize}
